\documentclass[10pt,a4paper]{amsart}
\usepackage[margin=19mm]{geometry}
\usepackage{amsmath,amssymb,mathtools}
\usepackage[scr=rsfs,cal=euler]{mathalfa}
\usepackage{xfrac}
\usepackage[unicode,hidelinks,bookmarks=false]{hyperref}
\newcommand{\ie}{i.e.\ }
\newcommand{\cf}{cf.\ }
\newcommand{\resp}{respectively\ }
\newcommand{\Subsection}{\subsection}
\providecommand{\nobreakdash}{\nobreak\hbox{-}}
\setlength{\emergencystretch}{2em}
\multlinegap0pt
\usepackage{amssymb}
\usepackage{mathtools}
\usepackage[shortlabels]{enumitem}
\usepackage{xcolor}
\newtheorem{lemma}{Lemma}
\newcommand{\A}{\mathbb{A}}
\newcommand{\B}{\mathbb{B}}
\newcommand{\C}{\mathbb{C}}
\newcommand{\D}{\mathbb{D}}
\newcommand{\E}{\mathbb{E}}
\newcommand{\F}{\mathbb{F}}
\newcommand{\G}{\mathbb{G}}
\renewcommand{\H}{\mathbb{H}}
\newcommand{\I}{\mathbb{I}}
\renewcommand{\L}{\mathbb{L}}
\newcommand{\N}{\mathbb{N}} % naturali
\renewcommand{\P}{\mathbb{P}}
\newcommand{\Q}{\mathbb{Q}}
\newcommand{\R}{\mathbb{R}} % reali
\title{Unified synthetic Ricci curvature lower bounds for Riemannian and sub-Riemannian structures}
\author{Davide Barilari, Andrea Mondino, Luca Rizzi}
\date{Source: arXiv:2211.07762v3 (CC BY 4.0)}
\begin{document}
\maketitle

\noindent\emph{Source and licence.} Unified synthetic Ricci curvature lower bounds for Riemannian and sub-Riemannian structures. Version arXiv:2211.07762v3 (CC BY 4.0), distributed by arXiv under the Creative Commons Attribution 4.0 International licence, http://creativecommons.org/licenses/by/4.0/, as recorded in the arXiv licence metadata for this exact version and shown on the abstract page https://arxiv.org/abs/2211.07762v3. Full licence notice: \texttt{licenses/arxiv-2211.07762-CC-BY-4.0.txt}.\par\smallskip

\noindent\textit{Editorial scope.} Counted unit: the article's lemma l:lemmaEcFc (source0.tex lines 3945--3958). The article's complete original proof of it is delivered with it (source0.tex lines 3959--4053, one \texttt{proof} environment of 95 lines). No mathematics is written, repaired or re-derived: the statement and the proof are the article's own lines, in the article's own notation, and no retained line is substituted or re-flowed.  Retained uncounted as the source-local context the proof needs: lines 3611--3613; lines 3646--3676; lines 3718--3724; lines 3757--3769; lines 3903--3916; lines 3919--3921.  Nothing was omitted from any retained span, so the package carries no omission record.  No retained line contains a citation, so the wrapper carries no bibliography and no bibliography entry.  No retained line contains a figure environment, an image-inclusion command or drawing code, so no image is delivered and the package declares no asset dependency.  Editorial formatting only, in addition: an AMS \texttt{amsart} preamble that restates the article's own declarations and macros used by the retained lines, namely the environments \texttt{lemma} on separate counters, together with the standard packages those lines need.  The retained environments print in this document as Lemma 1 in order of appearance, whereas the article prints the same items with the numbering of its own full document.  The complete native source of the article as distributed in the arXiv e-print for this exact version is bundled unchanged as \texttt{sources/133479/source0.tex}.
\begin{equation}\label{eq:sigmaframe}
\sigma(\partial_{h_i},\partial_{h_j})=0, \qquad \sigma(\partial_{h_i},\vec{h}_j)=\delta_{ij},\qquad \sigma(\vec{h}_i,\vec{h}_j) =  \sum_{\mu=1}^n h_\mu c_{ij}^\mu.
\end{equation}

\begin{lemma}[Fundamental computations]\label{l:fundamentalcomputations}
The following formulas hold on $T^*\mathcal{O}$:
\begin{align}
\dot{\partial}_{v_\mu} & = -u_\ell c_{\ell A}^\mu \partial_{h_A},\\
\dot{\partial}_{u_i}  & = - \vec{u}_i - u_\ell c_{\ell A}^i \partial_{h_A},\\
\dot{{\vec{h}}}_{A}   &  =  - h_B c_{A \ell}^B \vec{u}_\ell +u_\ell c_{\ell A}^B \vec{h}_B + u_\ell h_B f_{\ell A C}^B \partial_{h_C},
\end{align}
where for lower-case latin indices $i,j,\ell=1,\dots,k$, for greek ones $\mu,\nu,\sigma=k+1,\dots,n$, for upper-case latin ones $A,B,C=1,\dots,n$, repeated indices are understood as summed over their range, and we defined the following smooth functions on $\mathcal{O}$:
\begin{equation}
f_{\ell A C}^B\in C^\infty(\mathcal{O}),\qquad f_{\ell A C}^B := X_\ell(c_{AC}^B) - X_A(c_{\ell C}^B) - c_{\ell A}^D c_{DC}^B +c_{\ell D}^B c_{AC}^D-c_{AD}^B  c_{\ell C}^D.
\end{equation}
We rewrite these equations in a compact notation as follows:\begin{align}
\dot{\partial}_v & = \A \cdot \partial_u + \B \cdot \partial_v, \label{eq:partialvdot}\\
\dot{\partial}_u & = -\vec{u} + \C \cdot \partial_u + \D \cdot \partial_v, \label{eq:partialudot}\\
\dot{\vec{u}} & =  \E \cdot \vec{u} -\A^* \cdot \vec{v} + \F \cdot \partial_u + \G \cdot \partial_v, \label{eq:vecudot} \\ 
\dot{\vec{v}} & = \H \cdot \vec{u}  -\B^* \cdot \vec{v} + \I \cdot \partial_u + \L \cdot \partial_v, \label{eq:vecvdot}
\end{align}
where the dot denotes multiplication of tuples\footnote{Here and in the following, the notation $V = O\cdot W$, where $O$ is an $n\times m$ matrix, $V=(V_1,\dots,V_n)$ is a $n$-tuple of vector fields, and similarly $W=(W_1,\dots,W_m)$, with $n,m \in \N$, means that $V_i = \sum_{j=1}^m O_{ij} W_j$.}, and where we defined the following smooth matrix-valued maps:
\begin{align}
\A & : T^*\mathcal{O} \to \mathrm{M}(n-k, k), & \A_{\mu i} & := - u_\ell c_{\ell i}^\mu,  \\
\B & : T^*\mathcal{O} \to \mathrm{M}(n-k, n-k), & \B_{\mu \nu} & := - u_\ell c_{\ell\nu}^\mu,\\
\C & : T^*\mathcal{O} \to \mathrm{M}(k, k), & \C_{i j} & := - u_\ell c_{\ell j}^i, \\
\D & : T^*\mathcal{O} \to \mathrm{M}(k, n-k),& \D_{i \mu} & := - u_\ell c_{\ell \mu}^i, \\
\E & : T^*\mathcal{O} \to \mathrm{M}(k, k), & \E_{ij} & := -h_Bc_{ij}^B + u_\ell c_{\ell i}^j,   \\
\F & : T^*\mathcal{O} \to \mathrm{M}(k, k), & \F_{ij} & := u_\ell h_B f_{\ell i j}^B, \\
\G & : T^*\mathcal{O} \to \mathrm{M}(k, n-k),& \G_{i \mu} & := u_\ell h_B f_{\ell i \mu}^B, \\
\H & : T^*\mathcal{O} \to \mathrm{M}(n-k, k), & \H_{\mu i} & := - h_B c^{B}_{\mu i}+u_\ell c_{\ell \mu}^i,  \\
\I & : T^*\mathcal{O} \to \mathrm{M}(n-k, k), & \I_{\mu i} & :=  u_\ell h_B f_{\ell \mu i}^B, \\
\L & : T^*\mathcal{O} \to \mathrm{M}(n-k, n-k),& \L_{\mu \nu} & := u_\ell h_B f_{\ell \mu \nu}^B.
\end{align}
\end{lemma}

\begin{lemma}\label{l:property}
If $f:\Lambda \to V$ is $u$-pseudo-homogeneous of degree $d$, then $\vec{H}(f)$ is also $u$-pseudo-homogeneous of degree $d$. In particular, if $\Lambda\subset T^*\mathcal{O}$ is relatively compact domain, $j\in \N$, there exists $C=C_{j,\Lambda}$ such that
\begin{equation}
\|\vec{H}^{(j)}(f)(\lambda)\| \leq C H(\lambda)^{d/2}, \qquad \forall\, \lambda \in \Lambda_{\neq 0},
\end{equation}
where $\|\cdot\|$ denotes a fixed arbitrary norm on $V$.
\end{lemma}

\begin{lemma}\label{l:lemmaEa}
For any $\bar{\lambda}\in T^*M$ there exist a $u$-star-shaped neighborhood $\Lambda \subset T^*M$ of $\bar{\lambda}$, $T>0$, and a smooth map $\P: [0,T]\times \Lambda_{\neq 0} \to \mathrm{GL}(\R^{n-k})$, such that for all $\lambda \in \Lambda_{\neq 0}$ the tuple
\begin{equation}
E_a|_{\lambda_t} = \frac{1}{\sqrt{2H(\lambda)}}\P(t,\lambda) \cdot\partial_v|_{\lambda_t}, \qquad t\in[0,T],
\end{equation}
is (a choice of) the $E_a$-component of the canonical frame along $\lambda_t = e^{t\vec{H}}(\lambda)$.

Furthermore, the map $\P$ and all its time-derivatives are uniformly bounded, that is for all $j\in \N$ there exists a constant $C=C_j>0$ such that
\begin{equation}
\|\partial_t^{j}\P(t,\lambda)\| \leq C ,\qquad \forall\, t \in [0,T],\; \lambda \in \Lambda_{\neq 0},
\end{equation}
where $\|\cdot\|$ denotes a fixed matrix norm.
\end{lemma}

\begin{lemma}\label{l:lemmaEbFb}
For any $\bar{\lambda}\in T^*M$ there exist a $u$-star-shaped neighborhood $\Lambda \subset T^*M$ of $\bar{\lambda}$, $T>0$, and smooth maps $\Q_i: [0,T]\times  \Lambda_{\neq 0} \to \mathrm{M}(n_i,m_i)$, with $n_i,m_i \in \N$, $i=1,\dots,5$, such that for all $\lambda \in \Lambda_{\neq 0}$ the tuples
\begin{align}
E_b|_{\lambda_t} & = \Q_1(t,\lambda)\cdot \partial_u|_{\lambda_t}  + \frac{1}{\sqrt{2H(\lambda)}}\Q_2(t,\lambda) \cdot \partial_v|_{\lambda_t} , & t\in [0,T],\\
F_b|_{\lambda_t} & = \Q_3(t,\lambda)\cdot \vec{u} |_{\lambda_t} + \Q_4(t,\lambda)\cdot\partial_u|_{\lambda_t}  + \frac{1}{\sqrt{2H(\lambda)}}\Q_5(t,\lambda)\cdot\partial_v|_{\lambda_t} ,  & t\in [0,T],
\end{align}
are (a choice of) the $E_b$ and $F_b$-components of the canonical frame along $\lambda_t = e^{t\vec{H}}(\lambda)$.

Furthermore, the maps $\Q_i$ and all their time-derivatives are bounded, that is for all $j\in \N$ there exists a constant $C = C_j>0$ such that
\begin{equation}
\|\partial_t^{j}\Q_i(t,\lambda)\| \leq C ,\qquad \forall\, t \in [0,T],\; \lambda \in \Lambda_{\neq 0},\; i=1,\dots,5,
\end{equation}
where $\|\cdot\|$ denotes a fixed matrix norm.
\end{lemma}

\begin{equation}\label{eq:Ebexplicit}
E_b|_{\lambda_t}  = \frac{1}{\sqrt{2H}} \left[ \P(t,\lambda)\A(\lambda_t) \cdot \partial_u|_{\lambda_t} +  (\dot{\P}(t,\lambda)+ \P(t,\lambda) \B(\lambda_t) )\cdot \partial_v|_{\lambda_t}\right],
\end{equation}

\begin{lemma}\label{l:lemmaEcFc}
For any $\bar{\lambda}\in T^*M_{\neq 0}$ there exist a $u$-star-shaped neighborhood $\Lambda \subset T^*M$ of $\bar{\lambda}$, $T>0$, and smooth maps $\Q_i: [0,T]\times  \Lambda_{\neq 0} \to \mathrm{M}(n_i,m_i)$, with $n_i,m_i\in \N$, $i=6,\dots,10$, such that for all $\lambda \in \Lambda_{\neq 0}$ the tuples
\begin{align}
E_c|_{\lambda_t} & = \Q_6(t,\lambda) \cdot \partial_u|_{\lambda_t}  + \frac{1}{\sqrt{2H(\lambda)}}\Q_7(t,\lambda)\cdot \partial_v|_{\lambda_t} , & t\in [0,T],\\
F_c|_{\lambda_t} & =  \Q_8(t,\lambda) \cdot \vec{u}|_{\lambda_t}  + \Q_9(t,\lambda) \cdot \partial_u|_{\lambda_t}  + \frac{1}{\sqrt{2H(\lambda)}}\Q_{10}(t,\lambda)\cdot \partial_v|_{\lambda_t} , & t\in [0,T],
\end{align}
are (a choice of) the $E_c$ and $F_c$-components of the canonical frame along $\lambda_t = e^{t\vec{H}}(\lambda)$.

Furthermore, the maps $\Q_i$ and all their time-derivatives are bounded, that is for all $j\in \N$ there exists a constant $C = C_j>0$ such that
\begin{equation}
\|\partial_t^{j}\Q_i(t,\lambda)\| \leq C ,\qquad \forall\, t \in [0,T],\; \lambda \in \Lambda_{\neq 0},\; i=6,\dots,10,
\end{equation}
where $\|\cdot\|$ denotes a fixed matrix norm.
\end{lemma}

\begin{proof}
Let $\Lambda$ and $T>0$ as in Lemma \ref{l:lemmaEa} (cf.\@ first paragraph of its proof). Let $\lambda_t = e^{t\vec{H}}(\lambda)$, for $\lambda \in \Lambda_{\neq 0}$ and $t\in [0,T]$. The $(2k-n)$-tuple $E_c|_{\lambda_t}$ is determined, up to a constant orthogonal transformation, by the following conditions along $\lambda_t$:
\begin{itemize}
\item[(i)] $E_c$ is vertical (i.e.\@ $\pi_* E_c =0$);
\item[(ii)] Darboux frame conditions: $\mathbb{0}=\sigma(\dot{E}_b,E_c)$ and $\mathbb{0}=\sigma(\ddot{E}_b,E_c)$;
\item[(iii)] normalization condition: $\mathbb{1} = \sigma(\dot{E}_c,E_c)$;
\item[(iv)] isotropic condition: $\mathbb{0}=\sigma(\dot{E}_c,\dot{E}_c)$.
\end{itemize}
We observe that if $\bar{E}_c$ is a tuple which verifies (i), (ii) and (iii) along $\lambda_t$, then all these conditions will be also verified by
\begin{equation}
E_c|_{\lambda_t} : = \psi(t)\cdot \bar{E}_c|_{\lambda_t}, \qquad t\in [0,T],
\end{equation}
for any smooth map $t\mapsto \psi(t) \in \mathrm{O}(\R^{2k-n})$. Then, we find first smooth tuple $\bar{E}_c$ such that (i), (ii) and (iii) are verified, and subsequently we will fix $\psi$ via (iv). 

By (i), we must have
\begin{equation}\label{eq:ansatzEc}
\bar{E}_c|_{\lambda_t}= U(t)\cdot\partial_u|_{\lambda_t} + V(t)\cdot \partial_v|_{\lambda_t}, \qquad t\in [0,T],
\end{equation}
for some smooth $t\mapsto U(t) \in \mathrm{M}(2k-n,k)$ and $t\mapsto V(t)\in \mathrm{M}(2k-n,n-k)$. From this, and Lemma \ref{l:fundamentalcomputations}, we obtain:
\begin{align}
\dot{\bar{E}}_c &  = -U\cdot \vec{u}+(\dot{U} + U\C + V\A)\cdot \partial_u + (\dot{V} + U\D + V\B)\cdot\partial_v, \\
\ddot{\bar{E}}_c & =  -(2\dot{U} + U\E +  U\C + V\A)\cdot \vec{u} +U\A^* \cdot \vec{v} + \mathrm{span}\{\partial_u,\partial_v\}, 
\end{align}
Similarly from \eqref{eq:Ebexplicit} in the proof of Lemma \ref{l:lemmaEbFb}, and Lemma \ref{l:fundamentalcomputations} we obtain:
\begin{align}
\dot{E}_b & = \frac{1}{\sqrt{2H}}\P\A \cdot \vec{u} + \mathrm{span}\{\partial_u,\partial_v\},\\
\ddot{E}_b & = \mathbb{T}\cdot\vec{u} - \frac{1}{\sqrt{2H}} \P\A\A^* \cdot\vec{v} + \mathrm{span}\{\partial_u,\partial_v\},
\end{align}
where, for brevity of notation, we set $\mathbb{T} = \mathbb{T}(t,\lambda)$, as
\begin{equation}\label{eq:defT}
\mathbb{T}(t,\lambda):= - \frac{1}{\sqrt{2H}}\Big[2(\partial_t\P)\A + 3\P\dot{\A}+\P(\A\E+\A\C+\B\A) \Big],
\end{equation}
where the matrix maps $\A$, $\B$, $\C$, $\E$ in the r.h.s.\@ are evaluated along $\lambda_t = e^{t\vec{H}}(\lambda)$, $\P = \P(t,\lambda)$, for $\lambda \in \Lambda_{\neq 0}$ and $t\in [0,T]$. Notice that $\A$ is $u$-pseudo-homogeneous of degree $1$, together with $\dot{\A}$, by Lemma \ref{l:property}. Therefore we obtain that $\mathbb{T}$ is uniformly bounded with all its time-derivatives, that is for all $j\in \N$ there exists $C = C_j$ such that
\begin{equation}
\|\partial_t^{j}\mathbb{T}(t,\lambda)\|\leq C_j,\qquad \forall\, t\in [0,T],\; \lambda \in \Lambda_{\neq 0}.
\end{equation}

Using relations \eqref{eq:sigmaframe}, and the fact that $\P$ and $\A\A^*$ are invertible on $\Lambda_{\neq 0}$ (cf.\@ proof of Lemma \ref{l:lemmaEa}), conditions (ii) are equivalent to:
\begin{equation}\label{eq:pairofeqs}
\A U^* = 0, \qquad \text{and}\qquad V^* = \sqrt{2H}(\P\A\A^*)^{-1}\mathbb{T}U^*,
\end{equation}
along $\lambda_t$ and where $\P = \P(t,\lambda)$ and $\mathbb{T} = \mathbb{T}(t,\lambda)$.
The second equation of \eqref{eq:pairofeqs} determines $V(t)$ in terms of $U(t)$. The first one is an orthogonality relation between the rows of $\A(\lambda_t)$ and those of $U(t)$. Recall from Lemma \ref{l:fundamentalcomputations} that $\A : T^*\mathcal{O} \to \mathrm{M}(n-k,k)$ has constant rank equal to $n-k$ on $T^*\mathcal{O}_{\neq 0}$, while $U: [0,T]\to \mathrm{M}(2k-n,k)$ must have rank $2k-n$. By applying the Gram-Schmidt process to a local orthogonal complement to the rows of $\A(\bar{\lambda})$ we can find a neighborhood $\Lambda' \subset T^*\mathcal{O}$ of $\bar{\lambda}$ and a smooth map
\begin{equation}
\A^\perp : \Lambda'_{\neq 0} \to \mathrm{M}(2k-n,k).
\end{equation}
such that 
\begin{equation}\label{eq:propertiesA}
\A^\perp \A^* = \mathbb{0}, \qquad \A^\perp \A^{\perp *} = |u|^2 \mathbb{1}.
\end{equation}
Furthermore, since $\A$ is $u$-homogeneous of degree one, we can take $\A^\perp$ to be also $u$-homogeneous of degree $1$ and $\Lambda'$ to be $u$-star-shaped. The $u$-star-shaped neighborhood $\Lambda$ of $\bar{\lambda}$ built at the beginning of the proof of Lemma \ref{l:lemmaEa} was arbitrary. Up to restriction of $\Lambda$ to a smaller $u$-star-shaped set, and up to taking a smaller $T$ (and since the flow of $\vec{H}$ is smooth and preserves $|u|$), we can assume that for all $\lambda \in \Lambda_{\neq 0}$ and $t\in [0,T]$, the extremals $\lambda_t = e^{t\vec{H}}(\lambda)$ take value in a common bounded neighborhood contained in $\Lambda'_{\neq 0}$, so that $\A^\perp(\lambda_t)$ is well-defined. We then set:
\begin{align}\label{eq:UandV}
U(t) := \frac{1}{\sqrt{2H}}\A^{\perp}, \qquad V(t)  :=-\A^{\perp}\mathbb{T}^*(\P\A\A^*)^{-1*},
\end{align}
where the matrix maps $\A$, $\A^\perp$ in the r.h.s.\@ are evaluated along $\lambda_t = e^{t\vec{H}}(\lambda)$ and $\P = \P(t,\lambda)$, $\mathbb{T}=\mathbb{T}(t,\lambda)$ for $\lambda \in \Lambda_{\neq 0}$ and $t\in [0,T]$. Since $\A$ and $\A^\perp$ are $u$-homogeneous of degree $1$, $\P(t,\lambda)$ and $\mathbb{T}(t,\lambda)$ are uniformly bounded with all time-derivatives and $\lambda \in \Lambda_{\neq 0}$, and using Lemma \ref{l:property}, we observe that for all $j\in \N$ there exists a $C=C_j>0$ independent on $\lambda$ such that
\begin{equation}\label{eq:estimatesUV}
\|\partial_t^j U(t)\| \leq C, \qquad \|\partial_t V(t)\| \leq C H(\lambda)^{-1/2}, \qquad \forall\, t\in [0,T],\; \lambda \in \Lambda_{\neq 0}.
\end{equation}
Estimates \eqref{eq:estimatesUV} will be used later.

With the definition \eqref{eq:UandV}, conditions (i) and (ii) are fully verified, while (iii) corresponds to the following normalization:
\begin{equation}
\mathbb{1} = UU^* = \frac{1}{2H}\A^{\perp} \A^{\perp *},
\end{equation}
along $\lambda_t$, which is verified by our choice in the normalization of $\A^\perp$ in \eqref{eq:propertiesA}. This concludes the construction of the tuple $\bar{E}_c$.

Set hence $E_c|_{\lambda_t} = \psi(t)\cdot \bar{E}_c|_{\lambda_t}$, for $\psi(t) \in \mathrm{O}(\R^{2k-n})$. Condition (iv) yields an ODE that determines $\psi$, once the initial condition is fixed, for example $\phi(0) = \mathbb{1}$. We obtain
\begin{equation}\label{eq:ODEpsi}
\dot{\psi}(t) = \frac{1}{2}\psi(t)\sigma(\dot{\bar{E}}_c|_{\lambda_t},\dot{\bar{E}}_c|_{\lambda_t}), \qquad \psi(0)=\mathbb{1}.
\end{equation}
Since $(t,\lambda)\mapsto \sigma(\bar{E}_c|_{\lambda_t},\bar{E}_c|_{\lambda_t})$ is skew-symmetric and smooth, and $\mathrm{O}(\R^{2k-n})$ is compact, \eqref{eq:ODEpsi} will have a unique solution in $\mathrm{O}(\R^{2k-n})$ for all $t\in [0,T]$, and for any $\lambda \in \Lambda_{\neq 0}$. We define then the smooth map
\begin{equation}
\Psi : [0,T]\times \Lambda_{\neq 0} \to \mathrm{O}(\R^{2k-n}),
\end{equation}
such that $t\mapsto \Psi(t,\lambda)$ is the solution of \eqref{eq:ODEpsi} on $[0,T]$ corresponding to $\lambda_t = e^{t\vec{H}}(\lambda)$.

The matrix in the r.h.s.\@ of \eqref{eq:ODEpsi} is given more explicitly by\label{L: check this}
\begin{align}
\sigma(\dot{\bar{E}}_c|_{\lambda_t},\dot{\bar{E}}_c|_{\lambda_t}) & = U\sigma(\vec{u},\vec{u}) U^* -2\mathcal{A}\left[(\dot{U}+U\C+V\A)U^*\right].
\end{align}
where $\mathcal{A}(M)=\tfrac{1}{2}(M-M^*)$ denotes the skew-symmetric part of a matrix $M$, and $U$, $V$ are defined in \eqref{eq:UandV}. By the uniform estimates for $U(t)$ and $V(t)$ in \eqref{eq:estimatesUV}, and since $\A$ is $u$-homogeneous of degree $1$, we see that $\sigma(\dot{\bar{E}}_c|_{\lambda_t},\dot{\bar{E}}_c|_{\lambda_t})$ and all its time-derivatives remain uniformly bounded for all extremals $\lambda_t =e^{t\vec{H}}(\lambda)$ for $t\in [0,T]$ and $\lambda \in \Lambda_{\neq 0}$. By Gronwall's lemma we deduce an analogous uniform boundedness property of $\Psi$ and its time-derivatives, that is for all $j\in \N$ there exists $C=C_j >0$ such that
\begin{equation}
\|\partial_t^{j}\Psi(t,\lambda)\| \leq C, \qquad \forall\, t\in [0,T]\;  \lambda \in \Lambda_{\neq 0}.
\end{equation}

Finally, to obtain the statement of the lemma, we set then
\begin{align}
\Q_6(t,\lambda) & := \frac{1}{\sqrt{2H(\lambda)}}\Psi(t,\lambda) \A^\perp(\lambda_t), \\
\Q_7(t,\lambda) & :=  -\sqrt{2H(\lambda)}\Psi(t,\lambda)\A^\perp(\lambda_t)\mathbb{T}^*(t,\lambda)\left[\P(t,\lambda)\A(\lambda_t)\A(\lambda_t)^*\right]^{-1*}.
\end{align}
Using the fact that $\A^\perp$ is $u$-homogeneous of degree one, that $\Psi,\mathbb{T},\P$ are uniformly bounded on $[0,T]\times \Lambda_{\neq 0}$ with all their time-derivatives, and applying Lemma \ref{l:property}, we obtain the desired uniform bounds on $\Q_6,\Q_7$.

To prove the analogous statement for the tuple $F_c$, we recall that by the structural equations
$F_c|_{\lambda_t} = -\dot{E}_c|_{\lambda_t}$. Then, making use of the Leibniz rule, we obtain the corresponding statement of the lemma, for suitably defined matrix-valued maps $\Q_8,\Q_9,\Q_{10}$ defined on $[0,T]\times \Lambda_{\neq 0}$, satisfying the required uniform bounds.
\end{proof}

\end{document}
